\section{综合发现与研究启示}

\begin{reportbox}{综合判断}
结合分布、关键词、主题和作者网络四类证据, 当前数据呈现出一个清晰结构: 计算机科学和生物医学贡献主要样本规模, 材料科学样本较小但具备跨领域潜力; LLM、RAG、知识图谱等主题处于明显活跃期; 作者合作网络中既有稳定社区, 也存在需要消歧复核的高连接节点。
\end{reportbox}

\subsection{发现一: 近五年数据覆盖较均衡}

2022--2026 年每年均超过 30,000 条分析记录, 没有明显年份断点。该特点使年度趋势、融合术语动量和主题占比具备较好的比较基础。相比继续追求总量扩张, 当前更重要的是控制来源偏差和字段完整度差异。

\subsection{发现二: 热点主题呈现“平台化 + 场景化”双重特征}

LLM 和 RAG 代表平台型方法快速升温, prompt engineering、vision language model 和 knowledge graph 则体现出方法向具体科研场景迁移的趋势。生物医学和材料科学中的相关主题虽然规模不同, 但都能观察到方法迁移信号, 说明跨领域应用是后续值得关注的方向。

\subsection{发现三: 合作网络集中与分散并存}

作者网络既存在高连接度节点, 也存在多个领域内部社区。高连接作者不一定等于真实单个研究者影响力最高, 因为同名和缩写会带来聚合风险; 但这些节点适合作为人工复核入口, 帮助识别跨主题桥接者和稳定合作团簇。

\subsection{发现四: 材料科学需要归一化解读}

材料科学样本量低于计算机科学和生物医学, 直接比较绝对词频会低估部分材料方向主题。对于 battery、catalyst、materials discovery 等方向, 更合理的解释方式是观察领域内占比、近年动量和代表论文, 而不是只看全局排行。

\subsection{研究启示}

\begin{center}
\begin{tabularx}{\textwidth}{>{\bfseries}p{33mm}X}
\toprule
启示 & 数据依据 \\
\midrule
热点识别应使用多证据交叉 & 显式关键词、题名摘要术语、年度占比、近期动量和主题模型需要共同支持。 \\
跨领域机会主要出现在方法迁移处 & LLM、RAG、知识图谱等方法在不同领域中出现相似上升信号。 \\
合作网络适合发现候选桥接者 & 中心性和社区结构能提示跨团簇连接, 但需要作者消歧复核。 \\
全文覆盖决定深度分析边界 & 摘要级趋势可信度较高, 全文级细粒度解释仍需扩大正文覆盖。 \\
\bottomrule
\end{tabularx}
\end{center}

\clearpage
